%==============================================================================
% Appendix I: Reduction-order Artifact in the Spin-1 Sector
%==============================================================================
\section{Reduction-order Artifact in the Spin-1 Sector}
\label{app:I}

This appendix records a concrete example of a reduction-order artifact in a
sector with a rank-deficient field-space metric.  If one projects to a
complementary direction before eliminating the auxiliary direction, a quantity
that is not a finite generalised eigenvalue can appear to be a physical
eigenvalue.  The purpose is to clarify why the Schur reduction used in
Appendix~\ref{app:E} must retain the auxiliary coupling.

\subsection{Projected Quotient Before Auxiliary Elimination}

For each spin-1 axis, the raw variables $(\omega_k,\delta_k)$ give the witness
matrices
\begin{equation}
  H_2 =
  \begin{pmatrix}
    52.6379 & -631.6544 \\
    -631.6544 & 3789.9295
  \end{pmatrix},
  \qquad
  G_2 =
  \begin{pmatrix}
    236.870435 & -355.305547 \\
    -355.305547 & 532.958162
  \end{pmatrix}.
\end{equation}
The kinetic metric has the null direction
\begin{equation}
  n=\left(1,\frac{2}{3}\right)^T.
\end{equation}
A simple complementary direction is
\begin{equation}
  u=(-2,3)^T.
\end{equation}
The Rayleigh quotient along this direction is
\begin{equation}
  \lambda_\mathrm{proj}
  =\frac{u^T H_2u}{u^T G_2u}
  =\frac{41899.77}{10007.77}
  \simeq 4.1867>0.
\end{equation}
This quotient is a ratio of projected pre-elimination coefficients.  It is not,
by itself, the finite generalised eigenvalue of the rank-deficient pencil.

\subsection{Test Against the Generalised Eigenvalue Equation}

Because $G_2$ is rank-deficient and $G_2n\simeq 0$, a projected Rayleigh quotient
need not solve the degenerate pencil equation
\begin{equation}
  \det(H_2-\lambda G_2)=0.
\end{equation}
Indeed, substituting the projected value gives
\begin{equation}
  \det(H_2-\lambda_\mathrm{proj}G_2)\simeq -2.2\times 10^6\neq 0,
\end{equation}
and also
\begin{equation}
  (H_2-\lambda_\mathrm{proj}G_2)u\neq 0.
\end{equation}
Thus $u=(-2,3)$ is not a generalised eigenvector, and
$\lambda_\mathrm{proj}\simeq 4.1867$ is not a finite generalised eigenvalue.

\subsection{Missing Auxiliary Coupling}

The difference between the projected quotient and the Schur-reduced eigenvalue
comes from the Hessian cross term with the non-dynamical auxiliary direction:
\begin{equation}
  H_{YY}=n^TH_2n\simeq 894.845>0,
  \qquad
  H_{XY}=n^TH_2u
    =\frac{24081847}{3750}
    \simeq 6421.8\neq 0.
\end{equation}
Consequently the $n$ direction is not gauge or redundancy but a second-class
auxiliary direction.  It must be eliminated by its algebraic equation of motion
before a physical spin-1 eigenvalue is assigned.  A reduction that drops
$H_{XY}$ omits the auxiliary backreaction.

\subsection{Comparison with the Schur-reduced Result}

The correct order of reduction is:
\begin{enumerate}
  \item identify the null direction $Y_k$ of the rank-deficient kinetic metric;
  \item verify that $Y_k$ is a second-class auxiliary direction;
  \item solve the algebraic equation of motion for $Y$ in the quadratic form
        including the $X$--$Y$ cross term;
  \item evaluate the physical finite eigenvalue using the Schur complement
        \begin{equation}
          H_\mathrm{eff}=H_{XX}-\frac{H_{XY}^2}{H_{YY}}.
        \end{equation}
\end{enumerate}
This gives
\begin{equation}
  \lambda_\phys
  =\frac{H_{XX}-H_{XY}^2/H_{YY}}{G_{XX}}
  \simeq -0.4183<0.
\end{equation}
The same Schur-reduced structure in the symbolic Ricci normalisation gives
\begin{equation}
  \lambda_\phys^\mathrm{symbolic}
  =-\frac{56L^2}{3r^2}<0.
\end{equation}
Thus the positive projected quotient before auxiliary elimination and the finite
Schur-reduced eigenvalue have different signs and different physical
classifications.  The correct classification is a tachyonic non-ghost mode.

\subsection{Basis-independence Lesson}

Using the complementary direction $u=(-2,3)$ is not the essential problem.  The
problem is identifying a pre-elimination projected quotient with the physical
generalised eigenvalue in a rank-deficient sector.  Starting from the same
complementary coordinate still gives a negative finite eigenvalue once the
auxiliary direction is eliminated correctly.

The practical lesson is therefore:
\begin{quote}
In a sector with a rank-deficient field-space metric, the naive Rayleigh quotient
on a complementary direction must not be identified with a physical generalised
eigenvalue.  If a second-class auxiliary direction is present, one must use
Schur elimination, or equivalently the finite eigenvalue of the degenerate
pencil after retaining the auxiliary coupling.
\end{quote}

\subsection{Relation to Reproducibility Checks}

In the first version of this manuscript the pre-Schur projected quotient was
identified with the physical generalised eigenvalue.  In the present version the
finite eigenvalue after Schur elimination is used instead.  The affected
verification targets are:

\begin{table}[H]
\centering
\begin{tabularx}{\textwidth}{X l}
\toprule
Verification target & Section/appendix \\
\midrule
The value $4.1867>0$ is a projected quotient and not a finite generalised
eigenvalue. & Appendix~\ref{app:I} \\
After Schur elimination, the finite generalised eigenvalue is
$\lambda_\phys\simeq -0.4183<0$ in the witness normalisation and
$-56L^2/(3r^2)<0$ in the symbolic normalisation. &
\S\ref{sec:spin1}, Appendix~\ref{app:E} \\
The corrected spin-1 mode is a tachyonic non-ghost mode with positive kinetic
norm. & \S\ref{sec:spin1}, Appendix~\ref{app:E} \\
The spin-1 instability is a curvature-scale tachyonic non-ghost mode, distinct
from frame-degeneracy endpoint divergences and radial decompactification. &
\S\ref{sec:spin1}, \S\ref{sec:discussion} \\
\bottomrule
\end{tabularx}
\caption{Reproducibility checks associated with the spin-1 reduction-order
artifact.}
\label{tab:appI_checks}
\end{table}

\subsection{Computational Verification}

\begin{itemize}
  \item Reproduces the pre-Schur projected quotient 
        and verifies that it is not a finite generalised eigenvalue:
        \texttt{paper03/spin1\_appI\_projection\_artifact.py}
\end{itemize}
